% §262 `print_cs` renders a control-sequence name in three shapes, and the
% trailing space is part of the rendering rather than an afterthought:
%
%     procedure print_cs(@!p:integer); {prints a purported control sequence}
%     begin if p<hash_base then {single character}
%       if p>=single_base then
%         if p=null_cs then
%           begin print_esc("csname"); print_esc("endcsname"); print_char(" ");
%           end
%         else  begin print_esc(p-single_base);
%           if cat_code(p-single_base)=letter then print_char(" ");
%           end
%       else if p<active_base then print_esc("IMPOSSIBLE.")
%       else print(p-active_base)
%     ...
%     else  begin print_esc(text(p));
%       print_char(" ");
%       end;
%
% So: an *active character* prints bare — no escape character and no trailing
% space. A *single-character* control sequence prints with the escape and gets
% a trailing space only when that character's catcode is currently `letter` —
% a live `\catcode` test, not a question about the alphabet. A control *word*
% always gets the space.
%
% `sprint_cs` is the same minus the error checks and minus every trailing
% space; it is what `\string` and `\csname` echoing use.
%
% print_cs backs `show_token_list` (§294), and so `\tracingmacros`,
% `\tracingcommands`, `\message`, `\write` (§1370 `write_out` calls
% `token_show`), `\meaning` and `\detokenize`.
\input plain
% An active character: bare, no escape, no space. Braces mark the edges so a
% dropped or added space is visible.
\message{<str=[\string~]>}
\message{<detok=[\detokenize{~}]>}
\def\usesactive{a~b}
\message{<meaning=[\meaning\usesactive]>}
% A control word: escape plus a trailing space, always.
\message{<word=[\detokenize{\hbox}]>}
\message{<wordstr=[\string\hbox]>}
% A single-character control sequence whose character is a letter: space.
\message{<oneletter=[\detokenize{\a}]>}
% …and one whose character is not: no space. (Not `\+` — plain makes it
% `\outer`, and `\detokenize` refuses to scan one.)
\message{<onepunct=[\detokenize{\,}]>}
\message{<onespace=[\detokenize{\ }]>}
% The letter test reads the *current* catcode. With `@` made a letter, `\@`
% becomes a control word for printing purposes and picks up the space.
\catcode`\@=11
\message{<atletter=[\detokenize{\@}]>}
\catcode`\@=12
\message{<atother=[\detokenize{\@}]>}
% The converse: a letter demoted to `other` loses the space. It has to be a
% letter that no control word still ahead of us is spelled with — demoting `a`
% takes `\catcode` and `\message` with it, and `z` is inside `\detokenize`.
\message{<qletter=[\detokenize{\q}]>}
{\catcode`\q=12 \message{<qdemoted=[\detokenize{\q}]>}}
\message{<qrestored=[\detokenize{\q}]>}
% `\string` never adds a space, whatever the catcode.
\catcode`\@=11
\message{<stratletter=[\string\@]>}
\catcode`\@=12
% `\csname\endcsname` — the null control sequence — prints both halves and a
% space.
\expandafter\message\expandafter{\expandafter<\expandafter n\expandafter u\expandafter l\expandafter l\expandafter c\expandafter s\expandafter =\expandafter [\expandafter\detokenize\expandafter{\csname\endcsname}]>}
% `\escapechar` chooses the escape character, and a value outside 0..255
% suppresses it entirely (§63).
\escapechar=`\!
\message{<escbang=[\detokenize{\hbox}]>}
\escapechar=-1
\message{<escnone=[\detokenize{\hbox}]>}
\escapechar=`\\
% An active character is unaffected by `\escapechar`: it never had an escape.
\escapechar=`\!
\message{<escactive=[\string~]>}
\escapechar=`\\
\end
